% !TEX TS-program = pdflatex
% !TEX encoding = UTF-8 Unicode

% This file is a template using the "beamer" package to create slides for a talk or presentation
% - Giving a talk on some subject.
% - The talk is between 15min and 45min long.
% - Style is ornate.

% MODIFIED by Jonathan Kew, 2008-07-06
% The header comments and encoding in this file were modified for inclusion with TeXworks.
% The content is otherwise unchanged from the original distributed with the beamer package.

\documentclass{beamer}
\setbeamercovered{invisible}



% Copyright 2004 by Till Tantau <tantau@users.sourceforge.net>.
%
% In principle, this file can be redistributed and/or modified under
% the terms of the GNU Public License, version 2.
%
% However, this file is supposed to be a template to be modified
% for your own needs. For this reason, if you use this file as a
% template and not specifically distribute it as part of a another
% package/program, I grant the extra permission to freely copy and
% modify this file as you see fit and even to delete this copyright
% notice. 


\mode<presentation>
{
  \usetheme{Warsaw}
  % or ...

  % or whatever (possibly just delete it)
}


\usepackage[english]{babel}
% or whatever

\usepackage[utf8]{inputenc}
\usepackage{qrcode}
\usepackage{hyperref}
\usepackage{ytableau}
\usepackage{longtable}
% or whatever

\usepackage{times}
\usepackage[T1]{fontenc}
\usepackage{tikz}
% Or whatever. Note that the encoding and the font should match. If T1
% does not look nice, try deleting the line with the fontenc.

%%% MATH RELATED

\title{MATH 180 Strategies of Problem Solving} 
\subtitle
{} % (optional)

\author[W.R. Casper] % (optional, use only with lots of authors)
{W.R. Casper}
% - Use the \inst{?} command only if the authors have different
%   affiliation.

\institute[California State University Fullerton] % (optional, but mostly needed)
{
  Department of Mathematics\\
  California State University Fullerton}
% - Use the \inst command only if there are several affiliations.
% - Keep it simple, no one is interested in your street address.

\subject{Talks}
% This is only inserted into the PDF information catalog. Can be left
% out. 



% If you have a file called "university-logo-filename.xxx", where xxx
% is a graphic format that can be processed by latex or pdflatex,
% resp., then you can add a logo as follows:

% \pgfdeclareimage[height=0.5cm]{university-logo}{university-logo-filename}
% \logo{\pgfuseimage{university-logo}}



% Delete this, if you do not want the table of contents to pop up at
% the beginning of each subsection:
\AtBeginSubsection[]
{
  \begin{frame}<beamer>{Outline}
    \tableofcontents[currentsection,currentsubsection]
  \end{frame}
}


% If you wish to uncover everything in a step-wise fashion, uncomment
% the following command: 

%\beamerdefaultoverlayspecification{<+->}


\begin{document}

\begin{frame}
  \titlepage
\end{frame}

\begin{frame}{Outline}
  \tableofcontents
  % You might wish to add the option [pausesections]
\end{frame}

% Since this a solution template for a generic talk, very little can
% be said about how it should be structured. However, the talk length
% of between 15min and 45min and the theme suggest that you stick to
% the following rules:  

% - Exactly two or three sections (other than the summary).
% - At *most* three subsections per section.
% - Talk about 30s to 2min per frame. So there should be between about
%   15 and 30 frames, all told.

\section{SoPS Lecture 14}

\begin{frame}{Review}
  \begin{block}{Last time}
    \begin{itemize}
\pause
      \item Connect 4
\pause
      \item weak solutions vs strong solutions
\pause
      \item visualizing a Connect 4 weak solution
    \end{itemize}
  \end{block}
\end{frame}

\subsection{Connect 4 Problems}

\begin{frame}{Combinatorics problem}
\pause
  \begin{block}{Problem:}
\pause
    You have $10$ identical marbles, a blue cup, and a red cup.\\
\pause
    How many different ways can we put the marbles into the cups?
  \end{block}
\end{frame}

\begin{frame}{Combinatorics problem}
\pause
\textbf{Solution:}\\
  \begin{itemize}
\pause
    \item $10$ in blue, $0$ in red;
\pause
    \item $9$ in blue, $1$ in red;
\pause
    \item $8$ in blue, $2$ in red;
\pause
    \item $\dots$
\pause
    \item $1$ in blue, $9$ in red;
\pause
    \item $0$ in blue, $10$ in red.
  \end{itemize}
\pause
$$\text{{\color{blue}11 total ways}}$$
\end{frame}

\begin{frame}{Combinatorics problem -- version 2}
\pause
  \begin{block}{Problem:}
\pause
    You have $10$ identical marbles, a blue cup, a red cup, and a yellow cup.\\
\pause
    How many different ways can we put the marbles into the cups?
  \end{block}
\end{frame}

\begin{frame}{Combinatorics problem -- version 2}
\pause
\textbf{Solution:}\\
  \begin{itemize}
\pause
    \item $10$ in blue, $0$ in red, $0$ in yellow;
\pause
    \item $9$ in blue, $1$ in red, $0$ in yellow;
\pause
    \item $8$ in blue, $2$ in red, $0$ in yellow;
\pause
    \item $\dots$
\pause
    \item $1$ in blue, $9$ in red, $0$ in yellow;
\pause
    \item $0$ in blue, $10$ in red, $0$ in yellow;
  \end{itemize}
\pause
$$\text{{\color{red}11 total ways}}$$
\end{frame}

\begin{frame}{Combinatorics problem -- version 2}
\textbf{Solution (continued):}\\
  \begin{itemize}
\pause
    \item $9$ in blue, $0$ in red, $1$ in yellow;
\pause
    \item $8$ in blue, $1$ in red, $1$ in yellow;
\pause
    \item $7$ in blue, $2$ in red, $1$ in yellow;
\pause
    \item $\dots$
\pause
    \item $1$ in blue, $8$ in red, $1$ in yellow;
\pause
    \item $0$ in blue, $9$ in red, $1$ in yellow;
  \end{itemize}
\pause
$$\text{{\color{red}10 total ways}}$$
\end{frame}

\begin{frame}{Combinatorics problem -- version 2}
\textbf{Solution (continued):}\\
  \begin{itemize}
\pause
    \item $8$ in blue, $0$ in red, $2$ in yellow;
\pause
    \item $7$ in blue, $1$ in red, $2$ in yellow;
\pause
    \item $6$ in blue, $2$ in red, $2$ in yellow;
\pause
    \item $\dots$
\pause
    \item $1$ in blue, $7$ in red, $2$ in yellow;
\pause
    \item $0$ in blue, $8$ in red, $2$ in yellow;
  \end{itemize}
\pause
$$\text{{\color{red}9 total ways}}$$
\end{frame}

\begin{frame}{Combinatorics problem -- version 2}
\textbf{Solution (continued):}\\
\pause
The pattern continues ... until:
  \begin{itemize}
\pause
    \item $0$ in blue, $0$ in red, $10$ in yellow;
  \end{itemize}
\end{frame}

\begin{frame}{Combinatorics problem -- version 2}
\textbf{Solution (continued):}\\
\pause
Now we can get a {\color{blue}Grand Total:} $11+10+9 +\dots + 1$.\\
\pause
Remember Gauss's formula
$$1 + 2 + 3 + \dots + n = \frac{n(n+1)}{2}.$$
\pause
Using this we get our final answer:
$$\text{{\color{blue}66 total ways}}$$
\end{frame}

\begin{frame}{Combinatorics problem -- version 3}
\pause
  \begin{block}{Problem:}
\pause
    You have $10$ identical marbles, a blue cup, a red cup, a yellow cup, and a green cup.\\
\pause
    How many different ways can we put the marbles into the cups?
  \end{block}
  \begin{itemize}
\pause
    \item we could approach it like before
\pause
    \item might need \emph{sum of squares formula}
\pause
    \item is there an {\color{red}easier way}?
  \end{itemize}
\end{frame}

\begin{frame}{Combinatorics problem -- general version}
  \begin{block}{Problem:}
    You have $m$ identical marbles and $n$ different cups.\\
    How many different ways can we put the marbles into the cups?
  \end{block}
\pause
  \begin{block}{Solution:}
    We will use an argument called \textbf{stars and bars}.
  \end{block}
\end{frame}

\begin{frame}{Stars and bars}
  \begin{itemize}
\pause
    \item  Label the cups $1$ through $n$.
\pause
    \item  Marble arrangments are described by $m$ stars with $n-1$ bars.
\pause
    \item  How many stars before the first bar is how many marbles in the first cup.
\pause
    \item  How many stars after the first bar and before the second is how many marbles in the second cup.
\pause
    \item  Continue this pattern until, the number of stars after the last bar repesents the number of marbles in the $n$'th cup.
  \end{itemize}
\end{frame}

\begin{frame}{Stars and bars -- example}
  $$\star\star\vert \star \vert\vert \star \star \star \vert \star \star$$
  \begin{itemize}
\pause
    \item  $4$ bars = $5$ cups\\
\pause
    \item  $8$ stars = $8$ marbles\\
\pause
    \item  arrangement:
    \begin{itemize}
      \item  two marbles in cup $1$;
      \item one in cup $2$;
      \item zero in cup $3$;
      \item three in cup $4$;
      \item two in cup $5$.
    \end{itemize}
  \end{itemize}
\end{frame}

\begin{frame}{Challenge:}
  \begin{block}{Problem:}
    Consider the following \textbf{stars and bars} expression:
    $$\vert\vert\star\star\star\vert\star\star\vert\star\star\star\vert$$
    \begin{itemize}
\pause
      \item  how many cups in this arrangement
\pause
      \item  how many marbles in this arrangement
\pause
      \item  describe how many marbles in each cup
    \end{itemize}
  \end{block}
\end{frame}

\begin{frame}{Challenge:}
  \begin{block}{Problem:}
    Suppose you have $5$ cups with $3$ marbles in the first cup $2$ in the second, $4$ in the third, $0$ in the fourth and $1$ in the fifth.
    Determine the stars and bars expression which describes this arrangement.
  \end{block}
\end{frame}

\begin{frame}{Stars and bars formula}
  \begin{block}{Theorem (Stars and bars formula):}
    The number of ways of putting $n$ indistinguishable marbles into $m$ different cups is $\binom{m+n-1}{m}$.
  \end{block}
\pause
  \begin{block}{Proof:}
    It's the same as the number of stars and bars expressions with $m$ stars and $n-1$ bars.  Each expression is a string with $m+n-1$ characters, of which you choose $m$ spots for stars and fill the rest with bars.  Thus there's $m+n-1$ choose $m$ total ways.
  \end{block}
\end{frame}

\begin{frame}{Stars and bars formula}
  Examples:
  \begin{itemize}
\pause
    \item with $10$ marbles and two cups: $\binom{10+2-1}{10} = \binom{11}{10} = 11$
\pause
    \item with $10$ marbles and three cups: $\binom{10+3-1}{10} = \binom{12}{10} = 66$
\pause
    \item with $10$ marbles and four cups: $\binom{10+4-1}{10} = \binom{13}{10} = 286$
\pause
    \item with $15$ marbles and seven cups: $\binom{15+7-1}{15} = \binom{21}{15} = 54264$
  \end{itemize}
\end{frame}

\begin{frame}{In-class homework:}
  \begin{block}{Activity:}
    Work on current in-class problems.
  \end{block}
\end{frame}

\subsection{Ready Set Bet}

\begin{frame}{Play the game}
  \begin{block}{Activity}
    \begin{itemize}
\pause
      \item play a couple rounds of Ready Set Bet
\pause
      \item what strategies seem to be the best?
    \end{itemize}
  \end{block}
\end{frame}

\subsection{Wrapping up}

\begin{frame}{Final thoughts}
  \begin{block}{Reflection}
    \begin{itemize}
\pause
      \item Explain in your own words how an expression with stars and bars represents an arrangement of identical marbles in different containers.
    \end{itemize}
  \end{block}
\begin{center}
Submit your responses by \textbf{Tonight!}
\end{center}
\end{frame}


\end{document}


